% SPDX-License-Identifier: CC-BY-4.0
\documentclass[11pt,letterpaper]{article}

\usepackage[T1]{fontenc}
\usepackage{amsmath,amssymb}
\usepackage{tikz}
\usepackage{pgfplots}
\pgfplotsset{compat=1.18}
\usepackage{algorithm}
\usepackage{algpseudocode}
\usepackage{booktabs}
\usepackage{enumitem}
\usepackage[hidelinks]{hyperref}

\newcommand{\N}{\mathbb{N}}
\newcommand{\Z}{\mathbb{Z}}
\newcommand{\Q}{\mathbb{Q}}
\newcommand{\R}{\mathbb{R}}
\newcommand{\C}{\mathbb{C}}
\newcommand{\Map}[3]{#1 \colon #2 \to #3}
\newcommand{\Mod}[1]{#1\text{-}\mathrm{Mod}}
\newcommand{\GroupHom}{\mathrm{Hom}_{\mathbf{Grp}}}
\newcommand{\RingHom}{\mathrm{Hom}_{\mathbf{Ring}}}

\newtheorem{theorem}{Theorem}
\newtheorem{lemma}[theorem]{Lemma}
\newtheorem{corollary}[theorem]{Corollary}
\newtheorem{definition}[theorem]{Definition}
\newtheorem{remark}[theorem]{Remark}
\newenvironment{proof}{\noindent\textit{Proof.}}{\hfill\(\square\)}

\title{Generalized Collatz Library}
\author{}
\date{\today}

\begin{document}
\maketitle
\section*{Gcd strata of a positive offset}

\section*{Gcd strata of the positive offset map}

Throughout, fix an offset $c\in\N$ and define the positive offset map
\[
 C(n)=\begin{cases}n/2,&2\mid n,\\ 3n+c,&2\nmid n.\end{cases}
\]
Division has precedence, and $C^{k}$ denotes $k$ original instructions, so
$C^{0}$ is the identity and $C^{k+1}=C\circ C^{k}$. All statements are over
$\N$; the state $0$ is even and fixed. For $m\in\N$ the \emph{stratum} of a
state $n$ is the gcd $\gcd(n,m)$.

The results below separate four obligations: the one-step gcd algebra under
$m\mid c$ and $\gcd(m,6)=1$; the explicit odd short cycle of the offset map;
the counting bound for one stratum; and the fate of positive seeds when the
offset is even. Nothing below is a statement about the original Collatz map,
and no statement pretends to classify periodic states other than the three
listed in Theorem~\ref{Thm:OffsetCycle}.

\section{The map and its strata}

\begin{lemma}[Gcd under division]\label{Lem:OffsetTwoMulGcd}
\leavevmode\par
If $\gcd(m,2)=1$, then for every $k\in\N$,
\[
 \gcd(2k,m)=\gcd(k,m).
\]
\end{lemma}

\begin{proof}
Write $d=\gcd(2k,m)$ and $e=\gcd(k,m)$. Since $d\mid m$ and $d\mid 2k$, and
since $\gcd(m,2)=1$, the divisor $d$ of $m$ is coprime to $2$. Hence
$d\mid k$: indeed $d\mid 2k$ and $d$ is coprime to $2$, so $d$ divides the
other factor. Therefore $d\mid k$ and $d\mid m$, so $d\mid e$.

Conversely $e\mid k$ and $e\mid m$, hence $e\mid 2k$ and $e\mid m$, so
$e\mid d$. By antisymmetry of divisibility, $d=e$.
\end{proof}

\begin{theorem}[One-step gcd preservation]\label{Thm:OffsetStepGcd}
\leavevmode\par
Let $m\mid c$ and $\gcd(m,6)=1$. For every $n\in\N$,
\[
 \gcd(C(n),m)=\gcd(n,m).
\]
\end{theorem}

\begin{proof}
Since $\gcd(m,6)=1$, the modulus $m$ is coprime to $2$ and to $3$, as both
divide $6$.

\emph{Even case.} If $2\mid n$, write $n=2k$. Then $C(n)=k$, because division
has precedence and applies first, and
$\gcd(C(n),m)=\gcd(k,m)=\gcd(2k,m)=\gcd(n,m)$ by Lemma~\ref{Lem:OffsetTwoMulGcd}.

\emph{Odd case.} If $2\nmid n$, then $C(n)=3n+c$. Put $d=\gcd(3n+c,m)$ and
$e=\gcd(n,m)$. We prove $d\mid e$ and $e\mid d$.

For $d\mid e$: the divisor $d$ satisfies $d\mid m$ and $d\mid 3n+c$. From
$d\mid m$ and $m\mid c$ we get $d\mid c$. Then $d\mid(3n+c)-c=3n$. Since
$d\mid m$ and $\gcd(m,3)=1$, the divisor $d$ is coprime to $3$, so $d\mid n$.
Together with $d\mid m$ this gives $d\mid e$.

For $e\mid d$: from $e\mid n$ we get $e\mid 3n$, and from $e\mid m$ and
$m\mid c$ we get $e\mid c$. Hence $e\mid 3n+c$, and with $e\mid m$ this gives
$e\mid d$.

Antisymmetry of divisibility yields $d=e$, that is,
\[
\gcd(C(n),m)=\gcd(n,m).
\]
\end{proof}

\begin{corollary}[Orbits stay in their stratum]\label{Cor:OffsetOrbitStratum}
Let $m\mid c$ and $\gcd(m,6)=1$. For every $n\in\N$ and every $k\in\N$,
\[
 \gcd(C^{k}(n),m)=\gcd(n,m).
\]
\end{corollary}

\begin{proof}
Induction on $k$. For $k=0$ both sides are $\gcd(n,m)$. For the step, apply
Theorem~\ref{Thm:OffsetStepGcd} to the state $C^{k}(n)$ and use the induction
hypothesis:
\[
 \gcd(C^{k+1}(n),m)=\gcd(C(C^{k}(n)),m)=\gcd(C^{k}(n),m)=\gcd(n,m).
\]
\end{proof}

\begin{theorem}[Iterating one step later]\label{Thm:OffsetIterateSuccFirst}
For every $k\in\N$ and every $n\in\N$,
\[
 C^{k+1}(n)=C^{k}(C(n)).
\]
\end{theorem}

\begin{proof}
Write $k+1=1+k$ and apply the iterate addition identity
$C^{a+b}=C^{b}\circ C^{a}$, which is proved by induction on $b$ from the
definition $C^{k+1}=C\circ C^{k}$. With $a=1$ and $b=k$ this gives
$C^{k+1}=C^{k}\circ C$, which is the claim.
\end{proof}

\section{The odd short cycle}

\begin{theorem}[Odd offset cycle]\label{Thm:OffsetCycle}
Let $2\nmid c$. Then
\[
 C(c)=4c,\qquad C(4c)=2c,\qquad C(2c)=c,
\]
and consequently $C^{3}(c)=c$.
\end{theorem}

\begin{proof}
Since $c$ is odd, $3c+c=4c$ is even; hence $C(c)=4c$, because $c$ is odd and
the affine branch applies. The number $4c$ is even, so $C(4c)=4c/2=2c$. The
number $2c$ is even, so $C(2c)=2c/2=c$. Composing the three identities gives
$C^{3}(c)=c$.
\end{proof}

\begin{remark}\label{Rem:OffsetCycleScope}
The theorem exhibits one periodic orbit of length three. It does not assert
that every periodic state lies on this orbit, nor that the orbit attracts any
particular set of seeds. The counting result below bounds how many states can
lie in one stratum containing the cycle; it is not an estimate of a basin
beyond that stratum.
\end{remark}

\begin{theorem}[Coprime-seed exclusion]\label{Thm:OffsetCoprimeSeedAvoidsCycle}
\leavevmode\par
Let $m\mid c$, $\gcd(m,6)=1$, $m>1$, and let $n$ satisfy $\gcd(n,m)=1$. Then
for every $k\in\N$,
\[
 C^{k}(n)\ne c,\qquad C^{k}(n)\ne 4c,\qquad C^{k}(n)\ne 2c.
\]
\end{theorem}

\begin{proof}
By Corollary~\ref{Cor:OffsetOrbitStratum} and $\gcd(n,m)=1$ we have
$\gcd(C^{k}(n),m)=1$ for every $k$. Suppose $C^{k}(n)$ equalled one of $c$,
$4c$, $2c$. Since $m\mid c$, each of $c$, $4c$, and $2c$ is a multiple of
$m$; hence $m\mid C^{k}(n)$. Then $m$ divides both $C^{k}(n)$ and $m$, so
$m\mid\gcd(C^{k}(n),m)=1$. A divisor of $1$ is at most $1$, contradicting
$m>1$. Therefore none of the three states is reached.
\end{proof}

\section{Counting one stratum}

\begin{lemma}[Multiples in an initial segment]\label{Lem:OffsetMultipleCountLe}
For all $m\ge1$ and all $N\in\N$,
\[
 \#\{0\le n<N:m\mid n\}\le \frac{N}{m}+1,
\]
where the quotient is integer division.
\end{lemma}

The formal counterpart is stated for all $m,N\in\N$ with the integer quotient
of the formal system; at $m=0$ the estimate is immediate, since the set has at
most one element.

\begin{proof}
Let $S=\{0\le n<N:m\mid n\}$ and consider the map $\varphi(n)=n/m$ on $S$. If
$\varphi(a)=\varphi(b)$ for $a,b\in S$, then $a=m(a/m)=m(b/m)=b$ by exactness
of division on multiples, so $\varphi$ is injective.

For $n\in S$ we have $n\le N-1$, hence $n/m\le(N-1)/m\le N/m$, and therefore
$\varphi(n)<N/m+1$. Thus $\varphi$ maps $S$ into the range $\{0,\dots,N/m\}$,
whose cardinality is $N/m+1$. Injectivity gives $\#S\le N/m+1$.
\end{proof}

\begin{theorem}[Stratum density]\label{Thm:OffsetMultipleDensity}
For all $m,N\in\N$,
\[
 m\cdot\#\{0\le n<N:m\mid n\}\le N+m.
\]
\end{theorem}

\begin{proof}
For $m=0$ the left side is $0$, so the claim is $0\le N$. For $m\ge1$,
multiply the inequality of Lemma~\ref{Lem:OffsetMultipleCountLe} by $m$:
\[
 m\cdot\#S\le m\left(\frac{N}{m}+1\right)=m\left(\frac{N}{m}\right)+m\le N+m,
\]
where $m\lfloor N/m\rfloor\le N$ because $m\lfloor N/m\rfloor$ is the remainder
subtracted from $N$.
\end{proof}

\begin{remark}\label{Rem:OffsetDensityMeaning}
Dividing by $N$ shows that the multiples of $m$ have upper density at most
$1/m$. A periodic basin contained in the multiples of $m$ therefore has upper
density at most $1/m$; for $m>1$ such a basin omits every seed coprime to
$m$. The bound is a counting statement about the multiples of $m$, proved
independently of any dynamical estimate; no bound is claimed for the other
strata.
\end{remark}

\section{Even offsets}

\begin{theorem}[Odd states grow and stay odd]\label{Thm:OffsetEvenOffsetOddStep}
Let $2\mid c$ and $2\nmid n$. Then $2\nmid C(n)$ and $n<C(n)$.
\end{theorem}

\begin{proof}
Since $n$ is odd, the affine branch applies: $C(n)=3n+c$. If $2\mid C(n)$,
then from $2\mid c$ and $2\mid 3n+c$ we get $2\mid 3n$, and since $3$ is odd,
$2\mid n$, a contradiction. Hence $C(n)$ is odd. For the size, $n>0$ because
$n$ is odd, and $C(n)=3n+c\ge 3n>n$ since $c\ge0$.
\end{proof}

\begin{corollary}[Oddness persists]\label{Cor:OffsetEvenOffsetOddPersist}
Let $2\mid c$ and $2\nmid n$. For every $j\in\N$, the state $C^{j}(n)$ is odd.
\end{corollary}

\begin{proof}
Induction on $j$. At $j=0$ the state is $n$, which is odd. For the step,
Theorem~\ref{Thm:OffsetEvenOffsetOddStep} applied to the odd state
$C^{j}(n)$ shows that $C^{j+1}(n)=C(C^{j}(n))$ is odd.
\end{proof}

\begin{theorem}[Linear growth of odd orbits]\label{Thm:OffsetEvenOffsetGrowthLower}
Let $2\mid c$ and $2\nmid n$. For every $j\in\N$,
\[
 n+j\le C^{j}(n).
\]
\end{theorem}

\begin{proof}
Induction on $j$. At $j=0$ both sides are $n$. For the step, Corollary
\ref{Cor:OffsetEvenOffsetOddPersist} makes $C^{j}(n)$ odd, so
Theorem~\ref{Thm:OffsetEvenOffsetOddStep} gives
$C^{j}(n)<C^{j+1}(n)$. The induction hypothesis $n+j\le C^{j}(n)$ then yields
$n+j+1\le C^{j+1}(n)$, as required.
\end{proof}

\begin{lemma}[Exact dyadic prefix]\label{Lem:OffsetHalvePrefix}
For all $c,u,v\in\N$,
\[
 C^{v}\left(2^{v}u\right)=u .
\]
\end{lemma}

\begin{proof}
Induction on $v$. At $v=0$ both sides are $u$. For the step, $2^{v+1}u$ is
even, so $C(2^{v+1}u)=2^{v}u$; by the induction hypothesis,
$C^{v}(2^{v}u)=u$. Composing,
$C^{v+1}(2^{v+1}u)=C^{v}(C(2^{v+1}u))=C^{v}(2^{v}u)=u$.
\end{proof}

\begin{theorem}[Even-offset unboundedness]\label{Thm:OffsetEvenOffsetUnbounded}
\leavevmode\par
Let $2\mid c$ and $n>0$. For every bound $B\in\N$ there is $k\in\N$ with
\[
 B<C^{k}(n).
\]
\end{theorem}

\begin{proof}
Write $n=2^{v}u$ with $u$ odd, which is possible because $n>0$. By
Lemma~\ref{Lem:OffsetHalvePrefix}, $C^{v}(n)=u$. Let $k=v+(B+1)$. The iterate
addition identity gives
\[
 C^{k}(n)=C^{v+(B+1)}(n)=C^{B+1}\left(C^{v}(n)\right)=C^{B+1}(u).
\]
Since $u$ is odd, Theorem~\ref{Thm:OffsetEvenOffsetGrowthLower} with $j=B+1$
gives $u+(B+1)\le C^{B+1}(u)$. Finally $B<u+B+1$, so $B<C^{k}(n)$.
\end{proof}

\section{Scaling by a power of three}

\begin{theorem}[Offset scaling]\label{Thm:OffsetStepScaling}
For all $a,n\in\N$, writing $C_{c}$ for the map with offset $c$,
\[
 C_{3^{a}}\left(3^{a}n\right)=3^{a}\,C_{1}(n).
\]
\end{theorem}

\begin{proof}
The number $3^{a}$ is odd, so $2\mid 3^{a}n$ if and only if $2\mid n$.

If $2\mid n$, both $n$ and $3^{a}n$ are even, so the division branch applies on
both sides, and exactness of division on the multiple $3^{a}n$ gives
\[
 C_{3^{a}}\left(3^{a}n\right)=\frac{3^{a}n}{2}=3^{a}\frac{n}{2}
 =3^{a}C_{1}(n),
\]
because $n/2=C_{1}(n)$.

If $2\nmid n$, then $2\nmid 3^{a}n$, so the affine branch applies on both
sides:
\[
 C_{3^{a}}\left(3^{a}n\right)=3\cdot3^{a}n+3^{a}=3^{a}(3n+1)=3^{a}C_{1}(n),
\]
where $3n+1=C_{1}(n)$ because $n$ is odd.
\end{proof}

\begin{corollary}[Iterated scaling]\label{Cor:OffsetIterateScaling}
For all $a,n,k\in\N$,
\[
 C_{3^{a}}^{k}\left(3^{a}n\right)=3^{a}\,C_{1}^{k}(n).
\]
\end{corollary}

\begin{proof}
Induction on $k$. At $k=0$ both sides are $3^{a}n$. For the step, apply
Theorem~\ref{Thm:OffsetStepScaling} to the state $C_{1}^{k}(n)$ and the
induction hypothesis:
\[
 C_{3^{a}}^{k+1}\left(3^{a}n\right)
 =C_{3^{a}}\left(C_{3^{a}}^{k}\left(3^{a}n\right)\right)
 =C_{3^{a}}\left(3^{a}C_{1}^{k}(n)\right)
 =3^{a}C_{1}^{k+1}(n).
\]
\end{proof}

\begin{remark}[Scope of the scaling identity]\label{Rem:OffsetScalingScope}
The corollary identifies the orbits of the offsets $1$ and $3^{a}$ on scaled
seeds; it does not identify the two maps, and it does not transport
periodicity between unscaled and scaled seeds without the scaling relation on
both sides. In particular it is not used to transfer the cycle of
Theorem~\ref{Thm:OffsetCycle} between different offsets.
\end{remark}

\section{Controls and excluded cases}

The hypotheses of Theorem~\ref{Thm:OffsetStepGcd} cannot be dropped.

\begin{enumerate}[label=(\roman*)]
 \item $c=1$, $m=5$: the modulus does not divide the offset, and
 $\gcd(C(3),5)=\gcd(10,5)=5\ne1=\gcd(3,5)$.
 \item $c=3$, $m=3$: the modulus divides the offset but
 $\gcd(m,6)=3\ne1$, and $\gcd(C(1),3)=\gcd(6,3)=3\ne1=\gcd(1,3)$.
 \item $c$ even, seed $1$: Theorem~\ref{Thm:OffsetEvenOffsetUnbounded}
 applies, and the orbit is unbounded, while for odd $c$ the seed $c$ is
 periodic by Theorem~\ref{Thm:OffsetCycle}.
 \item $c=0$: the map is $n\mapsto n/2$ on even and $n\mapsto3n$ on odd
 inputs, the homogeneous map $T_{3,2}$ of the multiplicatively independent
 family. Every divisor of $0$ is admissible in
 Theorem~\ref{Thm:OffsetStepGcd}, so the numbered results apply at $c=0$ with a
 vacuous gcd identity.
\end{enumerate}

\section{Formal correspondence}

Every statement above is formalized in \texttt{lean/OffsetObstruction.lean}.
The correspondence is listed in Table~\ref{Table:OffsetObstructionLean}.

\begin{table}[htbp]
\centering
\begin{tabular}{ll}
\toprule
Statement & Lean declaration \\
\midrule
Lemma~\ref{Lem:OffsetTwoMulGcd} & \texttt{OffsetTwoMulGcd} \\
Theorem~\ref{Thm:OffsetStepGcd} & \texttt{OffsetStepGcd} \\
Corollary~\ref{Cor:OffsetOrbitStratum} & \texttt{OffsetOrbitStratum} \\
Theorem~\ref{Thm:OffsetCycle} & \texttt{OffsetCycleStepSelf} \\
 & \texttt{OffsetCycleStepFour} \\
 & \texttt{OffsetCycleStepTwo} \\
 & \texttt{OffsetCycle} \\
Theorem~\ref{Thm:OffsetCoprimeSeedAvoidsCycle} & \texttt{OffsetCoprimeSeedAvoidsCycle} \\
Lemma~\ref{Lem:OffsetMultipleCountLe} & \texttt{OffsetMultipleCountLe} \\
Theorem~\ref{Thm:OffsetMultipleDensity} & \texttt{OffsetMultipleDensity} \\
Theorem~\ref{Thm:OffsetEvenOffsetOddStep} & \texttt{OffsetEvenOffsetOddStep} \\
Corollary~\ref{Cor:OffsetEvenOffsetOddPersist} & \texttt{OffsetEvenOffsetOddPersist} \\
Theorem~\ref{Thm:OffsetEvenOffsetGrowthLower} & \texttt{OffsetEvenOffsetGrowthLower} \\
Lemma~\ref{Lem:OffsetHalvePrefix} & \texttt{OffsetHalvePrefix} \\
Theorem~\ref{Thm:OffsetEvenOffsetUnbounded} & \texttt{OffsetEvenOffsetUnbounded} \\
Theorem~\ref{Thm:OffsetStepScaling} & \texttt{OffsetStepScaling} \\
Corollary~\ref{Cor:OffsetIterateScaling} & \texttt{OffsetIterateScaling} \\
Theorem~\ref{Thm:OffsetIterateSuccFirst} & \texttt{OffsetIterateSuccFirst} \\
\bottomrule
\end{tabular}
\caption{Correspondence between the numbered statements and the Lean module.}
\label{Table:OffsetObstructionLean}
\end{table}

The finite verifier in \texttt{go/OffsetObstruction.go} checks the one-step
gcd identity, orbit strata, the cycle, the avoidance statement, the counting
bounds, the dyadic prefix, the growth statement, and the scaling identity on
bounded machine-word ranges, with negative controls for the hypotheses. Those
checks test the implementation; the universal statements are the theorems
above.

\end{document}
